% !TEX root = ../notes.tex

\documentclass[../notes.tex]{subfiles}

\begin{document}

\section{April 22}
Here we go.

\subsection{The Honda Formal Group}
Last time, we discussed a machine which takes in a pair $(R_0,\Gamma_0)$ of a classical $\FF_p$-algebra $R_0$ and a $p$-divisible group $\Gamma_0$, and it outputs some weakly even periodic $\mathbb E_\infty$-ring $E$ which knows about deformations.
\begin{example}
	If $R_0$ is a perfect field $k$ and $\Gamma_0$ is some connected $p$-divisible group of height $n$, then one finds that $\pi_*E(k,\Gamma_0)$ is $W(k)[[v_1,\ldots,v_{n-1}]][u,1/u]$. Morally, upon taking $\op{Spf}$, we are asking for $\{v_1,\ldots,v_{n-1}\}$ to go to nilpotents, and $u$ goes to a root of $v_n$, which are the coefficients of $x^{p^n}$ in the series $[p](x)$.
\end{example}
\begin{example}[Honda]
	Suppose $k=\FF_{p^n}$ and $\Gamma_0$ is the Honda formal group of height $n$, meaning that $\Gamma_0$ comes from the formal group law $F(x,y)\in\FF_p[[x,y]]$ for which $[p](x)=x^{p^n}$. (This uniquely specifies the formal group law.)
\end{example}
There is some content in taking $n$ twice here, which is seen in the endomorphisms.
\begin{definition}[endomorphism]
	An \textit{endomorphism} of a formal group law $F$ is a power series $\mc C$ for which
	\[\mc C(F(x,y))=F(\mc C(x),\mc C(y)).\]
\end{definition}
\begin{example}
	For the Honda formal group $\Gamma_0$ over $\FF_{p^n}$, there is an endomorphism $S$ given by the power series $x^p$, which successfully preserves the formal group law. It turns out that the endomorphism ring is given by
	\[\frac{W(\FF_{p^n})\langle S\rangle}{\left(S^n-p,Sa-a^\sigma S\right)_{a\in W(\FF_{p^n})}}.\]
	Here, $W(\FF_{p^n})$ only gets an action
\end{example}
\begin{notation}[Morava stabilizer group]
	We let $\OO_n$ be the endomorphism ring of the Honda formal group law of height $n$ (over $\FF_{p^n}$). We may also set $\mathbb S_n\coloneqq\OO_n^\times$ and $\mathbb G_n\coloneqq\mathbb S_n\rtimes\op{Gal}(\FF_{p^n}/\FF_p)$, and we define $E_n$ be $E(\FF_{p^n},\Gamma_0)$, where $\Gamma_0$ is the underlying Honda formal group.
\end{notation}
\begin{example}
	For $n=1$, one has $\OO_1=\ZZ_p$ because $S$ is then just $p$, and the Frobenius acts trivially on $\FF_p$.
\end{example}
\begin{remark}
	In general, $\OO_n^\times$ is a profinite group.
\end{remark}
\begin{remark}
	The $\mathbb E_\infty$-ring $E_n$ functorially receives an action by $\mathbb G_n$. This action turns out to be suitably continuous; for example, it is continuous on $\pi_0$. Furthermore, it will turn out that $L_{K(n)}\mathbb S=E_n^{h\mathbb G_n}$, and we in fact have a Galois extension. For example, this will imply that the category of $K(n)$-local spectra is equivalent to the category of $K(n)$-local $E_n$-modules with continuous $\mathbb G_n$-action.
\end{remark}
\begin{example}
	For $n=1$, one has $E_1=\mathrm{KU}^\land_p$ (by comparing homotopy groups) and $\mathbb G_1=\mathbb S_1=\ZZ_p^\times$, so $L_{K(1)}\mathbb S=\mathrm{KU}_p^{h\ZZ_p^\times}$. So if $p>2$, we can find a topological generator $\ell\in\ZZ_p^\times$, and one finds that
	\[L_{K(1)}X=\op{eq}\left(L_{K(1)}(\mathrm{KU}_p\otimes X)\stackrel\ell\rightrightarrows L_{K(1)}(\mathrm{KU}_p\otimes X)\right).\]
	For $p=2$, one has $\ZZ_2^\times=\{\pm1\}\times(1+4\ZZ_2)$, so one should replace $\mathrm{KU}_p$ in the above formula with $\mathrm{KO}_p=\mathrm{KU}_p^{h\{\pm1\}}$.
\end{example}
Note that there is an obvious issue between large and small primes in the above example. Thus, in the sequel, we will content ourselves with focusing on the (less complicated) larger primes.

\subsection{Light as a Feather}
We've been mentioning some continuous actions for a while, so we should probably explain what we mean by this. For applications, we will mostly be interested in actions by profinite groups on light solid condensed spectra. In the modern day, we know how to make sense of this on light solid condensed spectra. Thus, we need to explain what the adjectives ``light solid'' do.
\begin{remark}
	Note that spectra contain the connective spectra, which in turn are contained in the category of $\mathbb E_\infty$-spaces. Here, recall that an $\mathbb E_\infty$-space can be described by iterating $\mathbb E_1$-spaces. Alternatively, we can describe them as product-preserving functors from the category of free $\mathbb E_\infty$-spaces generated by finite sets of points to the category of spaces.
\end{remark}
Now, Segal has provided an explicit description of the full subcategory of $\mathbb E_\infty$-spaces which comes from the finite collection of points.
\begin{definition}[spans]
	Define $\op{Span}(\mathrm{FinSets})$ to be the $2$-category whose objects are finite sets, its morphisms are correspondences of finite sets, and the compositions are given by composition of correspondences as follows: $X\from Y\to Z$ and $Z\from Y'\to W$ compose by taking the pullback correspondence $Y\from Y''\to Y'$ which fits into the diagram
	% https://q.uiver.app/#q=WzAsNixbMCwyLCJYIl0sWzEsMSwiWSJdLFsyLDIsIloiXSxbMywxLCJZJyJdLFs0LDIsIlciXSxbMiwwLCJZJyciXSxbNSwxXSxbMSwwXSxbMSwyXSxbMywyXSxbMyw0XSxbNSwzXSxbNSwyLCIiLDAseyJzdHlsZSI6eyJuYW1lIjoiY29ybmVyIn19XV0=&macro_url=https%3A%2F%2Fraw.githubusercontent.com%2FdFoiler%2Fnotes%2Fmaster%2Fnir.tex
	\[\begin{tikzcd}[cramped]
		&& {Y''} && \\
		& Y && {Y'} \\
		X && Z && W
		\arrow[from=1-3, to=2-2]
		\arrow[from=1-3, to=2-4]
		\arrow["\lrcorner"{anchor=center, pos=0.125, rotate=-45}, draw=none, from=1-3, to=3-3]
		\arrow[from=2-2, to=3-1]
		\arrow[from=2-2, to=3-3]
		\arrow[from=2-4, to=3-3]
		\arrow[from=2-4, to=3-5]
	\end{tikzcd}\]
	to give a correspondence $X\from Y''\to W$. The $2$-morphisms are natural isomorphisms of correspondences.
\end{definition}
\begin{remark}[Segal]
	It turns out that the full subcategory of $\mathbb E_\infty$-spaces given by the free $\mathbb E_\infty$-spaces on finite collections of points is exactly $\mathrm{Span}(\mathrm{FinSet})$.
\end{remark}
Thus, we can think about an $\mathbb E_\infty$-space as a product-preserving functor $\op{Span}(\mathrm{FinSet})\to\mathrm{Spaces}$.
\begin{example}
	The $\mathbb E_\infty$-space $\Omega^\infty X$ corresponds to the functor which sends $[n]$ to $(\Omega^\infty X)^n$. Then a map $A\to B$ of finite sets corresponds to the span $B\from A=A$, and it sends the maps $B\to\Omega^\infty X$ to the corresponding map
	\[A\to B\to\Omega^\infty X.\]
	To see that we have some notion of an operation, we note that the span $A=A\to B$ should be able to lift maps $A\to\Omega^\infty X$ to a map $B\to\Omega^\infty$, which should morally be given by ``summing over fibers,'' so we have some notion of summation!
\end{example}
\begin{remark}
	It is not too surprising that one can find a description of $\mathbb E_\infty$-spaces in terms of finite sets because $\mathrm{FinSet}$ is the free $\mathbb E_\infty$-space generated by a point.
\end{remark}
We can now define light solid condensed spectra all at once.
\begin{definition}[light solid condensed]
	A \textit{light solid condensed spectrum} $F$ is a product-preserving functor
	\[\op{Span}(\mathrm{CountableSet})\to\mathrm{Spaces},\]
	where $\pi_0F(\mathrm{pt})$ is a group. Here, the correspondences $A\from B\to C$ in $\op{Span}(\mathrm{CountableSets})$ is required to have the map $B\to C$ have finite fibers.
\end{definition}
\begin{remark}
	If $X$ is a countable set, then $F(X)$ should be thought of as the maps from $X$ to the light solid condensed spectrum whose image is compactly supported. Intuitively, such a thing should have some topology and some group operation.
\end{remark}
\begin{remark}
	This functor has found a way to encode both the topology and the group operation at the same time.
\end{remark}

\subsection{Sting like a Bee}
We now go back and discuss the adjectives ``light'' and ``solid'' and ``condensed'' separately. Here is the classical story.
\begin{definition}[light]
	A \textit{light profinite set} is a topological space which is a sequential limit (in $\mathrm{Top}$) of finite discrete spaces.
\end{definition}
\begin{remark}
	There is a forgetful functor which takes a light profinite set to the pro-category of finite sets, where we forget all the topological information.
\end{remark}
\begin{definition}[condensed]
	A \textit{light condensed abelian group} is a functor $A$ from the opposite category of light profinite sets to the category of abelian groups, which is a sheaf in the sense that it satisfies the following.
	\begin{listalph}
		\item $A(\emp)=0$.
		\item $A(X\sqcup Y)=A(X)\times A(Y)$.
		\item For any continuous surjection $X\onto Y$ of light profinite sets, one has $A(Y)=\lim(A(X)\rightrightarrows A(X\times_YX))$.
	\end{listalph}
\end{definition}
\begin{example}
	One can send any abelian group $A$ and give it the discrete topology, which allows us to define a light condensed abelian group: it should take a light profinite set $X$ to the abelian group of continuous maps $X\to A$.
\end{example}
We now want to make these definitions homotopy-coherent. Morally, we have seen that we are trying to build sheaves on the category of light profinite sets.
\begin{definition}[hypercover]
	Fix a simplicial light profinite set $S_\bullet$ is a \textit{hypercover}. Then a map $S_0\to S$ is a \textit{hypercover} if and only if the induced maps
	\[S_n\to S(\del\Delta^n),\]
	where $S(\del\Delta^n)$ is calculated in the category of light profinite sets.
\end{definition}
\begin{definition}[condensed]
	A \textit{light condensed spectrum} $X$ is a functor from the opposite category of light profinite sets to spectra such that $X(A\sqcup B)=X(A)\times X(B)$ and
	\[X(S)=\lim_{n\in\Delta}X(S_n)\]
	for any hypercover $S_\bullet\to S$.
\end{definition}
\begin{definition}[condensed]
	Fix a presentable category $\mc C$. Then we define $\mc C^\ell$ to be the category of \textit{light condensed} objects in $\mc C$ to be the category of hypersheaves on $\op{ProFin}^\ell$ valued in $\mc C$. Here, a hypersheaf is a functor which satisfies the sheaf condition for hypersheaves, and $\op{ProFin}^\ell$ is the category of light condensed sets.
\end{definition}
\begin{example}
	The constant sheaf associates a ``discrete'' condensed object to any object.
\end{example}
\begin{remark}
	One can also view $\mc C^\ell$ is the tensor product of $\mathrm{Spaces}^\ell\otimes\mc C$, where the tensor product takes place in the category of presentable categories.
\end{remark}
And now we can give the definition of solid. Here is the classical notion.
\begin{definition}[solid]
	A light condensed abelian group $A$ is \textit{solid} if and only if the following holds: for each light profinite set $S=\lim_nS_n$ and map $\ZZ[S]\to M$, there is a unique extension fitting in the following triangle.
	% https://q.uiver.app/#q=WzAsMyxbMCwwLCJcXFpaW1NdIl0sWzAsMSwiXFxkaXNwbGF5c3R5bGVcXGxpbV9uXFxaWltTX25dIl0sWzEsMCwiTSJdLFswLDJdLFswLDFdLFsxLDIsIiIsMix7InN0eWxlIjp7ImJvZHkiOnsibmFtZSI6ImRhc2hlZCJ9fX1dXQ==&macro_url=https%3A%2F%2Fraw.githubusercontent.com%2FdFoiler%2Fnotes%2Fmaster%2Fnir.tex
	\[\begin{tikzcd}[cramped]
		{\ZZ[S]} & M \\
		{\displaystyle\lim_n\ZZ[S_n]}
		\arrow[from=1-1, to=1-2]
		\arrow[from=1-1, to=2-1]
		\arrow[dashed, from=2-1, to=1-2]
	\end{tikzcd}\]
\end{definition}
\begin{remark}
	The limit $\lim_n\ZZ[S_n]$ turns out to not depend on the choice of presentation $S=\lim_nS_n$. Accordingly, we may write it as $\ZZ[S]^\square$.
\end{remark}
Roughly speaking, we are trying to ensure that $M$ is complete (namely, it can ``fill in'' from the limit). Now, here is the coherent definition.
\begin{definition}[solid]
	A light condensed spectrum $X$ is \textit{solid} if and only if $\pi_nX$ is solid for all $n$.
\end{definition}
\begin{remark}
	It turns out that the category of solid light condensed spectra (or abelian groups) is closed under colimits and limits.
\end{remark}
\begin{remark}
	Fix a light profinite set $S$. If $\{e_i\}_{i\in I}$ is the abelian group of continuous functions $S\to\ZZ$, then
	\[\ZZ[S]^\square=\prod_I\ZZ,\]
	where the product is taken in the category of light condensed abelian groups. It turns out that these objects are compact projective generators, so we can understand our functors purely on these objects. Unwinding this can eventually recover spans.
\end{remark}

\end{document}